\thispagestyle{empty}
\newpage
\cleardoublepage\phantomsection
\addcontentsline{toc}{chapter}{Abstract}\index{Abstract}

\begin{center}
{\Large \bf Abstract}\\
\end{center}
\vspace{10pt}

The “Research Paper Analysis and Critique” project aims to assist researchers, students, and academicians by providing an intelligent system that automates the process of analyzing, summarizing, and critiquing research papers. This report outlines the project’s objectives, methodology, and future prospects in advancing academic research support.

With the exponential rise in the number of published research papers across disciplines, it has become increasingly difficult for scholars to efficiently identify relevant studies, evaluate their quality, and validate the authenticity of claims. To address this challenge, “Research Paper Analysis and Critique” was developed as a web-based application that leverages cutting-edge AI models and natural language processing techniques. The system extracts text from uploaded research papers, summarizes the content, checks citation integrity, detects potential plagiarism, verifies factual accuracy, and provides a structured critique of methodology and conclusions.

The application employs transformer-based models for abstractive summarization, TF-IDF with cosine similarity for plagiarism detection, CrossRef and Semantic Scholar APIs for citation validation, and fact-checking services for claim verification. By combining these modules, the system offers users an all-in-one academic assistant capable of reducing the manual burden of reviewing large volumes of literature.

Beyond assisting individual users, “Research Paper Analysis and Critique” provides value to the wider academic community by promoting research quality and integrity. Insights derived from analysis can help institutions, publishers, and reviewers detect questionable claims, improve peer-review processes, and support evidence-based knowledge dissemination.

Overall, the system represents a significant step forward in integrating AI into academic research workflows. By simplifying paper evaluation, enhancing critical insights, and enabling trustworthy research validation, the application has the potential to transform how researchers engage with scientific literature in the future.

\vspace{0.5cm}
\noindent\textbf{Key Points:} research paper analysis, AI critique, summarization, citation validation, fact-checking, plagiarism detection.
